\section{Schlussfolgerung, Abgrenzung und weitere Forschung}
\label{sec:conclusion}

Die geplante Arbeit schließt eine wesentliche Lücke in der praktischen Anwendung von 
\gls{ids}-Standards. Während die technische Validierung von Modellen bereits etabliert ist, 
bietet die aktive Modellanreicherung mittels \gls{ids} ein erhebliches Potenzial zur 
Prozessoptimierung in der Angebotsphase.

\subsection{Zusammenfassung und erwarteter Beitrag}
Der angestrebte wissenschaftliche Beitrag dieser Arbeit liegt in dem Nachweis, dass 
die \gls{ids} als Datenquelle eingesetzt werden kann, um die Nachbearbeitung von 
\gls{ifc}-Modellen zu unterstützen und das dieser Prozess technisch möglich ist. 
Durch die Entwicklung und Evaluation des Prototyps wird aufgezeigt, wie die Datentiefe des Modells 
für die anschließende Mengenermittlung verbessert werden kann. 
Damit liefert die Arbeit sowohl ein technisches Artefakt als auch methodische Erkenntnisse.

\subsection{Abgrenzung der Forschung}
Um den Fokus der Untersuchung zu schärfen und den Umfang der Arbeit einzugrenzen, werden 
folgende Aspekte explizit abgegrenzt:

\begin{itemize}
    \item \textbf{Fachliche Abgrenzung:} \\
    Der Prototyp stellt die technische Infrastruktur zur Anreicherung bereit. 
    Die inhaltliche Verantwortung (z. B. die korrekte fachliche Zuweisung von Attributen 
    zu Bauteilen durch den Anwender) verbleibt beim Menschen. Eine bautechnische 
    Plausibilitätsprüfung der zugewiesenen Werte findet nicht statt.
    \item \textbf{Thematische Abgrenzung:} \\
    Die Arbeit konzentriert sich ausschließlich auf die Anreicherung von alphanumerischen 
    Informationen (Properties/Attributes). Eine Korrektur von Modellgeometrien wird nicht 
    betrachtet.
    \item \textbf{Vollständige Prozessautomatisierung:} \\
    Der Fokus dieser Arbeit liegt auf der technischen Schnittstelle zwischen 
    \gls{ids} und \gls{ifc}. Die Validität der anschließenden Mengenermittlung auf 
    Basis \gls{ids}-konformer Modelle wird als Zielprozess vorausgesetzt. 
    \item \textbf{Software:} \\
    Die Umsetzung erfolgt prototypisch auf Basis von Open-Source Technologien. 
    Eine Integration in kommerzielle Autorensoftware wird nicht durchgeführt.
    \item \textbf{Phasen- und Geographische Abgrenzung:} \\
    Die Untersuchung beschränkt sich auf die Angebotsphase im Hochbau unter Berücksichtigung 
    der im österreichischen Raum üblichen Anforderungen.
\end{itemize}

\subsection{Grenzen und weitere Forschung}
Eine Analyse wie Fachanwender:innen durch den Prozess der Modellanreicherung geführt werden können, 
ist im Rahmen dieser primär technisch orientierten Arbeit nicht vorgesehen und bietet Raum für zukünftige 
Usability-Studien zur Gestaltung intuitiver Benutzeroberflächen. 
Zukünftige Forschungsarbeiten könnten darauf aufbauen, indem KI-gestützte Ansätze zur automatisierten 
Zuordnung von \gls{ids}-Anforderungen integriert werden oder die Skalierbarkeit des Ansatzes auf 
weitere Lebenszyklusphasen oder den Tiefbau untersucht wird. Darüber hinaus könnten zukünftige Untersuchungen 
das Potenzial von IDS als operative Datengrundlage für weitere Anwendungsfälle analysieren. Beispielsweise 
könnte IDS nicht nur zur nachträglichen Anreicherung, sondern bereits als Modellierungshilfe eingesetzt 
werden, um Anwender:innen schon während der Geometrieerstellung aktiv bei der korrekten Attributierung zu unterstützen. 
Ergänzend bleibt festzuhalten, dass die Wirksamkeit des entwickelten Artefakts in der Praxis eng an die 
Verbreitung von \gls{ids}-basierten Workflows, der Qualität der Modellgeometrien und der in \gls{ids} 
definierten Anforderungen geknüpft ist.